\section{The exact third steering map with two evolving parents}
\label{tr:section}

The actual third transition supplies a nonvertical second phase and
a nonzero second vorticity. This section derives the original next
steering system from that state. Its prescribed horizontal component
determines all three phases explicitly. The resulting six older
equations retain both parent temperatures, both parent vorticities,
and their different physical diffusion rates. We prove an explicit
positive interval of the actual steering evolution and the exact
shooting and holding equations. A third return root is not asserted
here; its endpoint estimates are a subsequent calculation.

\subsection{Actual source data and prescribed horizontal component}

Use the parameters and actual endpoint $t_1$ of
Theorem~\ref{tt:transition}. In particular retain
\[
 \begin{gathered}
 a=\sin s_2,\quad b=r_b\sin s_3,\quad
 c=A_0\cos s_*,\quad C=A_0\sin s_*,\\
 s=s_3=\frac{L_3\sigma_1}{\sigma_2},\quad
 L=L_3,\quad \Gamma=\Gamma_3=\sigma_2\sin s.
 \end{gathered}
\]
The constants $a,b,c,C$ are positive. The quantity
$\nu=\sqrt{A_0}$ remains the source time scale, whereas $d>0$
is physical viscosity and thermal diffusivity. The original physical
frequencies remain $K_j=\mu\widehat\lambda_j$ with
$\mu=\lambda_0=\lceil\sqrt{A_0}/2\rceil$.
Put
\begin{equation}\label{tr:entry}
 Z_1=\zeta_{3,1}(t_1)>0,\qquad
 \chi_1=|\zeta_3(t_1)|>0,\qquad
 P_{3,1}=-\Theta_3(t_1)>0,\qquad
 v_1=\frac{\sigma_2\Omega_3(t_1)}{K_3\chi_1\Theta_3(t_1)}>0.
\end{equation}
All these numbers are the actual diffusive transition values; neither
$\chi_1$ nor $Z_1/\sin s$ is replaced by one.
Write $\tau=\Gamma(t-t_1)$, $T=1+L^{-1}$, and retain exactly
the source profiles $\eta,\eta_2$ and $H=\|\eta_2\|_\infty$.
Thus $\eta$ is a smooth nondecreasing step from zero to one on
$[0,1]$, and $\eta_2\in C_c^\infty((0,1))$ is nonnegative
with integral one. For the source trial parameter $m\in[0,3]$, set
\begin{equation}\label{tr:Z}
 Z_m(\tau)=
 \begin{cases}
 Z_1(1-\eta(\tau)),&0\leq\tau\leq1,\\
 -mL\sin s\,\chi_1\eta_2(L(\tau-1)),&1\leq\tau\leq T,\\
 0,&\tau\geq T.
 \end{cases}
\end{equation}
The source pulse parameter in (3.12) is $h_3=mL\sin s$;
its full range $[0,3L\sin s]$ is preserved. The multiplier
$\chi_1$ in the negative pulse is required by that equation.
Flatness of the step and interior support of the pulse make
\eqref{tr:Z} smooth at the two joins, with every derivative zero
there apart from its displayed constant value at $\tau=0$.
The prescribed family is independent of $m$ before $\tau=1$.
For every integer $r\geq1$ its physical derivatives on the two
pieces are exactly
\begin{equation}\label{tr:Z_derivatives}
 \partial_t^r Z_m=
 \begin{cases}
 -Z_1\Gamma^r\eta^{(r)}(\tau),&0\leq\tau\leq1,\\
 -m\chi_1\sin s\,L(\Gamma L)^r
          \eta_2^{(r)}(L(\tau-1)),&1\leq\tau\leq T.
 \end{cases}
\end{equation}
In particular the pulse derivative in the angular control contains
$\Gamma L^2$, with its full original coefficient.

\subsection{Complete physical coordinate equations and inverse}

Let $J=\left(\begin{smallmatrix}0&-1\\1&0\end{smallmatrix}\right)$.
Use the exact determinant coordinates of \eqref{tt:coordinates}:
\begin{equation}\label{tr:coordinates}
 \begin{gathered}
 D=h^2+a^2,\qquad
 x=\frac{hg-ab}{D},\qquad y=\frac{ag+hb}{D},\qquad
 R=x^2+y^2=\frac{g^2+b^2}{D},\\
 E_1=-K_1\Theta_1,\qquad E_2=-K_2\Theta_2,\qquad
 N_2=a(E_1+c)+Ch,\qquad
 N_3=(E_1+c)y+Cx+bE_2.
 \end{gathered}
\end{equation}
Here $D=|\zeta_2|^2$ and $R=|\zeta_3|^2$ are scalars.
For a prescribed $Z=Z_m(\Gamma(t-t_1))$ put
\begin{equation}\label{tr:inverse}
 \begin{gathered}
 Y=\sqrt{R-Z^2},\qquad \zeta_3=(Z,Y),\\
 \zeta_2=\left(\frac{gZ-bY}{R},\frac{gY+bZ}{R}\right),
 \qquad
 \zeta_1=\left(\frac{xZ-yY}{R},\frac{xY+yZ}{R}\right).
 \end{gathered}
\end{equation}
The square root specifies the original positive third second component.
On the component domain used below, $h,g,x,y>0$ and all three
second components are positive. The exact angular inverse is
\begin{equation}\label{tr:angles}
 \begin{gathered}
 \phi_3=\arcsin(Z/\sqrt R),\qquad
 \phi_2=\phi_3-\arctan(b/g),\\
 \phi_1=\phi_2-\arctan(a/h),\qquad
 \alpha=s_*-\phi_1.
 \end{gathered}
\end{equation}
Both inverse trigonometric functions take their principal real branches.
These formulas retain the original angular lift at $t_1$; continuity
in this positive-component domain precludes a change by $2\pi$.
Indeed the actual transition has
$\phi_3=\theta+\arctan(b/g)\in(-\pi/2,\pi/2)$, as follows
from \eqref{tt:revivalbounds}, $g>3/4$ and $b<4a$.
The principal inverse therefore agrees with the actual transition
lift, rather than selecting a new representative. An alternative
unwrapped initial lift would retain its fixed initial $2\pi$ multiple
throughout the same formulas.

The complete closed older system in physical time is
\begin{equation}\label{tr:closed}
 \begin{aligned}
 \dot h&=a\Omega_1,\\
 \dot g&=\Omega_1\frac{2ahg+(h^2-a^2)b}{D}+b\Omega_2,\\
 \dot E_1&=-C\Omega_1-dK_1^2E_1,\\
 \dot\Omega_1&=-E_1\frac{xZ-yY}{R}-dK_1^2\Omega_1,\\
 \dot E_2&=-\frac{N_2}{D}\Omega_2-dK_2^2DE_2,\\
 \dot\Omega_2&=-E_2\frac{gZ-bY}{R}-dK_2^2D\Omega_2.
 \end{aligned}
\end{equation}
Its six initial values are inherited from $t_1$ without resetting any
older amplitude. Neither $\Omega_3$ nor $\Theta_3$ enters this
older system, because the third shear does not enter its own parent.
The exact angular control reconstructed from it is
\begin{equation}\label{tr:control}
 \begin{gathered}
 \mathcal F_3=
 \frac{\Omega_1y\zeta_{1,1}+(b\Omega_2/D)\zeta_{2,1}}{Y},
 \qquad
 \dot\alpha=\mathcal F_3-\frac{\dot Z}{Y},\\
 \dot Z=\Gamma\,\partial_\tau Z_m.
 \end{gathered}
\end{equation}
Thus the prescribed component and the source control (3.12) are
exactly equivalent; no independent angular forcing was substituted.

The complete original deformation matrices are
\begin{equation}\label{tr:matrices}
 \begin{gathered}
 M_1=R_\alpha,\qquad
 M_2=R_{\alpha-s_*}
 \begin{pmatrix}1&-(h-\cos s_2)/a\\0&1\end{pmatrix},\\
 \mathcal L=(\zeta_2\ \zeta_3),\qquad
 M_3=\mathcal L(t)^{-T}\mathcal L(t_b)^T.
 \end{gathered}
\end{equation}
Here $t_b=t_2^b$ is the original third insertion time, not the
third steering-entry time. All three matrices retain their original
activation data, have determinant one, and obey their respective
physical transport equations.

\begin{proof}[Exact forward and inverse calculation]
The elementary identities
\[
 hx+ay=g,\qquad hy-ax=b,\qquad x^2+y^2=R
\]
follow by expanding \eqref{tr:coordinates}. The inverse
\eqref{tr:inverse} gives
\[
 |\zeta_1|^2=1,\quad |\zeta_2|^2=D,\quad
 |\zeta_3|^2=R,\quad
 \zeta_2=h\zeta_1-aJ\zeta_1,\quad
 \zeta_3=x\zeta_1-yJ\zeta_1.
\]
For example $\zeta_1=(xI+yJ)\zeta_3/R$, whose squared
length is $(x^2+y^2)|\zeta_3|^2/R^2=1$.
Also $\zeta_2=(gI+bJ)\zeta_3/R$; multiplying
$(hI-aJ)(xI+yJ)=(gI+bJ)$ proves the first relative formula.
Taking their oriented determinants gives
\[
 \det(\zeta_1,\zeta_2)=-a,\qquad
 \det(\zeta_2,\zeta_3)=-b.
\]
The dot products are respectively $h$ and $g$, so the construction
is inverse to the determinant coordinates, with the original signs.
The polar expressions for these same identities give
\eqref{tr:angles}, since multiplication by $J$ rotates the vector
$(\sin\phi,\cos\phi)$ in the direction of decreasing $\phi$.

Differentiating $x,y$ with the first two equations of
\eqref{tr:closed} gives the exact formulas
\begin{equation}\label{tr:xyR}
 \dot x=\Omega_1y+\frac{bh}{D}\Omega_2,
 \qquad \dot y=\frac{ba}{D}\Omega_2,
 \qquad \dot R=2\Omega_1xy+\frac{2bg}{D}\Omega_2.
\end{equation}
In particular the second parent's vorticity changes $y$; it has
not disappeared in this representation. Differentiating
\eqref{tr:angles} gives
\[
 \dot\alpha=-\frac{\dot Z}{Y}
       +\frac{Z\dot R}{2RY}
       -\frac{b\dot g}{g^2+b^2}-\frac{a\dot h}{D}.
\]
For an explicit check write $A=ax+hy$, so
$\dot g=\Omega_1A+b\Omega_2$. Then
\[
 bA+a^2R=(hy-ax)(ax+hy)+a^2(x^2+y^2)=Dy^2.
\]
The last two terms in the angular derivative are consequently
$-\Omega_1y^2/R-b^2\Omega_2/(DR)$.
Combining these with $Z\dot R/(2RY)$ gives exactly
$(\Omega_1y(xZ-yY)+(b\Omega_2/D)(gZ-bY))/(RY)$,
which equals
$(\Omega_1y\zeta_{1,1}+(b\Omega_2/D)\zeta_{2,1})/Y$.
This proves \eqref{tr:control}.

For a direct check of the phase dynamics define
\[
 D_{<2}=\dot\alpha J+
       \Omega_1J\zeta_1\otimes\zeta_1,\qquad
 D_{<3}=D_{<2}+\frac{\Omega_2}{D}
                         J\zeta_2\otimes\zeta_2.
\]
Since $\phi_1+\alpha=s_*$, one has
$\dot\zeta_1=\dot\alpha J\zeta_1$.
Differentiating the two relative formulas above using
\eqref{tr:xyR} gives, without dropping either shear,
\[
 \dot\zeta_2=\dot\alpha J\zeta_2+a\Omega_1\zeta_1
             =-D_{<2}^T\zeta_2,
 \qquad
 \dot\zeta_3=\dot\alpha J\zeta_3+\Omega_1y\zeta_1
                   +\frac{b\Omega_2}{D}\zeta_2
             =-D_{<3}^T\zeta_3.
\]
Their first third component is $\dot Z$ by
\eqref{tr:control}; their second is the derivative of the selected
square root. The older physical gradient is
\[
 G_{<3}=-(E_1+c+hE_2)\zeta_1-(C-aE_2)J\zeta_1.
\]
Hence $-J\zeta_2\cdot G_{<2}=N_2$ and
$-J\zeta_3\cdot G_{<3}=N_3$, proving the coefficients
$a_2=N_2/(K_2D)$ and $a_3=N_3/(K_3R)$.
The damped physical amplitude rows and
$\dot\Omega_j=K_j\zeta_{j,1}\Theta_j-dK_j^2|\zeta_j|^2\Omega_j$
are now exactly the last four equations of \eqref{tr:closed}.

The added second shear has transpose annihilating $\zeta_2$.
Thus the columns of $\mathcal L$ both solve the third transport
equation, and $\det\mathcal L=-b$. Differentiation of
$\mathcal L^{-T}\mathcal L(t_b)^T$ proves
$\dot M_3=D_{<3}M_3$ and its original value $I$ at $t_b$.
Differentiating the two earlier matrices proves their equations as
in \eqref{tt:matrices}; the constants $a$ and $\cos s_2$ are
their original data. Conversely, the physical phase equations give
the first two rows of \eqref{tr:closed} by differentiating the dot
products; the amplitude rows give the other four. These calculations
prove the full forward and inverse equivalence.
\end{proof}

The exact loss equations continue through the prescribed steering:
\begin{equation}\label{tr:losses}
 \dot N_2=-adK_1^2E_1,\qquad
 \dot N_3=-d(K_1^2E_1y+bK_2^2DE_2).
\end{equation}
Indeed differentiating $N_2$ cancels $-aC\Omega_1$ with
$Ca\Omega_1$. Differentiating $N_3$ gives the shear terms
\[
 -C\Omega_1y+C\Omega_1y
 +\frac{b\Omega_2}{D}
       \{a(E_1+c)+Ch-N_2\}=0;
\]
the remaining terms are precisely \eqref{tr:losses}.
The value of $\dot Z$, and therefore the possibly large derivative
of the compressed pulse, does not enter either cancellation.
An additional exact algebraic relation useful for continuing the
older temperature calculation is
\begin{equation}\label{tr:temperature_inverse}
 E_1=\frac{N_2-Ch}{a}-c,\qquad
 E_2=\frac{N_3}{b}-\frac{N_2y}{ab}+\frac Ca.
\end{equation}
To check the second relation insert the first in $N_3$ and use
$ax-hy=-b$. Thus one may equally use $N_2,N_3$ instead of
$E_1,E_2$ in the older system, with the displayed inverse and
the explicit diffusion losses, without assuming either numerator
constant at positive $d$.

\subsection{An explicit actual interval of smooth steering}

\begin{theorem}\label{tr:local}
For every original datum and positive physical diffusivity constructed
by \eqref{tge:Y3}--\eqref{tge:d3}, the actual third transition is
followed by a unique smooth solution of \eqref{tr:closed} and the
source steering control on
\[
 0\leq\tau\leq\tau_0:=\frac{\sin s}{8192},
 \qquad
 t_1\leq t\leq t_1+\frac1{8192\sigma_2}.
\]
All $m\in[0,3]$ give the same solution there. Every older
temperature is negative, all source phase second components are
positive, and the original matrices and their inverses are smooth.
On this interval the third temperature and vorticity remain strictly
negative and the exact source quotient below is positive.
\end{theorem}

\begin{proof}
This interval is proved from the actual entry margins, without
assuming bounds on a future trial. Use the dimensionless state
\[
 U=(h,g,E,F,\omega,V)
  =\left(h,g,\frac{E_1}{\sigma_1^2},
       \frac{E_2}{\sigma_2^2},
       \frac{\Omega_1}{\sigma_1},
       \frac{\Omega_2}{\sigma_2}\right).
\]
Define $\epsilon=\sigma_1/\Gamma\leq2/L$,
$q=\sigma_1/\sigma_2=s/L$, $\bar c=c/\sigma_1^2\leq1/4$,
$\bar C=C/\sigma_1^2\leq a/32$, and
$\delta_j=dK_j^2/\Gamma$. Primes in this proof denote $\tau$
derivatives. The exact scaled older equations are
\begin{equation}\label{tr:scaled}
 \begin{aligned}
 h'&=a\epsilon\omega,\\
 g'&=\epsilon\omega\frac{2ahg+(h^2-a^2)b}{D}
                          +\frac b{\sin s}V,\\
 E'&=-\epsilon\bar C\omega-\delta_1E,\qquad
 \omega'=-\epsilon E\zeta_{1,1}-\delta_1\omega,\\
 F'&=-q\epsilon\frac{a(E+\bar c)+\bar Ch}{D}V-\delta_2DF,
 \qquad
 V'=-\frac{F\zeta_{2,1}}{\sin s}-\delta_2DV.
 \end{aligned}
\end{equation}
Keep the closed cube $|U-U(t_1)|_\infty\leq1/128$.
The actual bounds \eqref{tt:olderbounds} and
\eqref{tt:revivalbounds} imply throughout this cube
\[
 \begin{gathered}
 63/64<h<4,\quad 11/16<g<6,\quad
 15/128<E<3,\quad 7/128<F<4,\\
 |\omega|<11,\quad |V|<1/16,\quad
 15/16<D<17.
 \end{gathered}
\]
Here the lower $h$ bound uses $a\leq1/512$ and
$h(t_1)\geq\sqrt{1-a^2}$. The bound on $V$ uses
$7040a/L<1/32$. For $0\leq\tau\leq1$,
\eqref{tr:Z} and \eqref{tt:newestcoefficientbox} give
\[
 0\leq Z\leq Z_1<\frac98\sin s<\frac1{128},
 \qquad b<4\sin s\leq4a.
\]
Consequently
\[
 R>\frac{(11/16)^2}{17}>\frac1{64},\qquad
 Y^2>\frac1{64}-\frac1{128^2}>\frac1{81},
 \qquad x>\frac1{32},\quad y>0.
\]
For $x$ use
$hg-ab>(63/64)(11/16)-4a^2>17/32$ and $D<17$.
Thus $Y>1/9$, $gY+bZ>0$, and $xY+yZ>0$ on this entire
known cube and time interval. The original positive-component
domain has a fixed positive margin; $R$ and $D$ never vanish.
All right-hand sides are smooth on a neighborhood of this cube.

Their sizes can be bounded before solving. The first four rows of
\eqref{tr:scaled} have absolute value less than one: use
$\epsilon\leq2/L$, $L\geq16384$, $\delta_1\leq1/128$,
$|\zeta_{1,1}|\leq1$, and
\[
 \left|\frac{2ahg+(h^2-a^2)b}{D}\right|
 \leq\frac{2ag}{h}+b<17a.
\]
The fifth row also has absolute value less than one, since
$a(E+\bar c)+\bar Ch<4a$, $q=s/L\leq a/L$,
and $\delta_2DF<68/128$.
Finally $|\zeta_{2,1}|\leq\sqrt D<5$ gives
\[
 |V'|<\frac{20}{\sin s}+\frac{17}{2048}
       <\frac{21}{\sin s}.
\]
Thus the norm of the complete vector field is strictly less than
$M=32/\sin s$ on this cube. Local existence follows from its
smooth integral equation. On any initial solution inside the cube,
$|U(\tau)-U(0)|_\infty<M\tau$. At $\tau=\tau_0$
this is at most $1/256$, strictly smaller than its radius $1/128$.
A first exit is impossible, and the same smooth compact-domain
bound extends the solution through $\tau_0$. The integral
equation and its bounded first derivatives prove uniqueness.
Since $\tau_0<1$, every trial has the same prescribed $Z$ there.
The reconstructed original matrices and control are smooth because
the denominators just bounded are positive.
They join the preceding transition smoothly: there
$\dot\zeta_{3,1}=(1-\eta_b)(\mathcal F_3-\mathcal F_2)\zeta_{3,2}$,
whose derivatives of every order vanish at $t_1$. The new prescribed
$Z$ has the same constant value and the same vanishing derivatives
there by \eqref{tr:Z_derivatives}. Both controls consequently have
the common endpoint jet of $\mathcal F_3$; differentiating their
smooth physical equations successively matches every state derivative.

On this actual interval $E_1,E_2,x,y>0$, so $N_3>0$.
The newest linear pair has nonnegative off-diagonal coefficients
$A=N_3/(K_3R)>0$ and $B=K_3Z\geq0$.
The three coefficients $A,B,\lambda=dK_3^2R$ are bounded on
the proved compact older trajectory. In the displayed original
pair coordinates, $C_*=\max(\lambda+A+B)$ bounds the induced
matrix $1$-norm. Its integral equation gives
$|\Theta_3(t)|+|\Omega_3(t)|\leq
(|\Theta_3(t_1)|+|\Omega_3(t_1)|)e^{C_*(t-t_1)}$
by Gronwall's inequality. This finite bound extends the unique
linear pair through the same closed interval.
Removing the common factor
$\exp(-dK_3^2\int R\,dt)$ gives a pair whose negative
temperature and negative vorticity coordinates have nonnegative
derivatives whenever both are positive. Their strictly positive
initial values and their integral equations therefore keep both
positive through this interval; a first zero contradicts that
integral representation. This proves the final sign assertion.
\end{proof}

\subsection{The source quotient, pulse equation, and exact return submanifold}

The newest physical pair associated to any constructed steering
geometry is
\begin{equation}\label{tr:physical_pair}
 \frac d{dt}\begin{pmatrix}\Theta_3\\\Omega_3\end{pmatrix}
 =\begin{pmatrix}
 -dK_3^2R&N_3/(K_3R)\\
 K_3Z&-dK_3^2R
 \end{pmatrix}
 \begin{pmatrix}\Theta_3\\\Omega_3\end{pmatrix}.
\end{equation}
For a fundamental matrix equal to $I$ at $t_1$, direct
differentiation of its two-column determinant gives
\begin{equation}\label{tr:wronskian}
 \det\mathcal U(t)=
 \exp\left(-2dK_3^2\int_{t_1}^tR(r)\,dr\right)>0.
\end{equation}
Indeed differentiating a determinant column by column multiplies it
by the trace of the displayed coefficient matrix; solving that scalar
equation gives \eqref{tr:wronskian}. A nonzero pair therefore never
becomes the zero vector. This does not by itself prevent an individual
temperature zero, which is why a finite quotient bound must still be
proved when extending through the full negative pulse.

In the chart $\Theta_3<0$ define the exact source variables
\begin{equation}\label{tr:quotient}
 \begin{gathered}
 v=\frac{\sigma_2\Omega_3}{K_3\chi_1\Theta_3},\qquad
 k=\frac{\chi_1N_3}{\sigma_2^2\sin s\,R},\qquad
 f_m=\frac{Z_m}{\chi_1\sin s},\qquad
 \delta_3=\frac{dK_3^2}{\Gamma},\\
 v'=f_m-kv^2,\qquad
 P_3=P_{3,1}\exp\left(\int_0^\tau(kv-\delta_3R)\,dr\right),\\
 \Theta_3=-P_3,\qquad
 \Omega_3=-\frac{K_3\chi_1}{\sigma_2}P_3v.
 \end{gathered}
\end{equation}
Differentiating the quotient cancels the two equal newest damping
terms and gives its displayed Riccati row. Differentiating the
exponential and product gives exactly \eqref{tr:physical_pair},
including both newest diffusion terms. Thus the map is an exact
inverse on every finite quotient interval; on the actual interval
of Theorem~\ref{tr:local} all its divisions and signs have already
been proved.

Put $z_1=Z_1/(\chi_1\sin s)$. The forcing is exactly
\begin{equation}\label{tr:forcing}
 f_m(\tau)=
 \begin{cases}
 z_1(1-\eta(\tau)),&0\leq\tau\leq1,\\
 -mL\eta_2(L(\tau-1)),&1\leq\tau\leq T,\\
 0,&\tau\geq T.
 \end{cases}
\end{equation}
In particular the source normalization removes $\chi_1$ from the
negative pulse exactly, while retaining it in $k$, in $z_1$, and in
the amplitude inverse. The actual transition estimates imply
\[
 \frac{63}{64}<\chi_1<\frac32,\qquad
 \frac13<v_1<\frac{64}{21},\qquad
 \frac5{16}<z_1<\frac87.
\]
The length bounds follow from \eqref{tt:coefficientvariation},
$1\leq R(t_a)\leq2$, and
$e^{-1/128}>63/64$, $2e^{1/64}<9/4$.
The bounds for $v_1$ use $1/2\leq\sigma_2\Omega_3/(K_3\Theta_3)\leq3$;
the bounds for $z_1$ use \eqref{tt:newestcoefficientbox}.
They are entry data, not estimates for the unconstructed full pulse.

The exact pulse integral equation, written on its quotient chart,
is
\begin{equation}\label{tr:pulse_integral}
 \begin{gathered}
 V_m(y)=v(1+y/L;m),\qquad
 H_2(y)=\int_0^y\eta_2(r)\,dr,\\
 V_m(y)=v(1;m)-mH_2(y)
       -\frac1L\int_0^y k(1+r/L;m)V_m(r)^2\,dr,
 \qquad 0\leq y\leq1.
 \end{gathered}
\end{equation}
It follows by substituting $\tau=1+y/L$ in the Riccati equation
and integrating, retaining the trial dependence of every older
coefficient. The required shooting endpoint is precisely $v(T;m)=0$,
because the temperature inverse is strictly negative on a finite
quotient interval. Formula \eqref{tr:pulse_integral} specifies that
endpoint without a frozen coefficient or a second-stage substitution.
Establishing its opposite endpoint signs and exclusion of temperature
zeros requires full-interval geometry and quotient estimates.
Those estimates have not been supplied by the local theorem.

The exact older forcing signs identify the additional analytic work.
For $Z\geq0$ in the positive $x,y,g,Y$ domain,
\[
 (gZ-bY)(gZ+bY)=R(DZ^2-b^2),\qquad
 (xZ-yY)(xZ+yY)=R(Z^2-y^2).
\]
Thus the second parent's horizontal phase changes sign precisely at
$Z=b/\sqrt D$, and the first parent's at $Z=y$.
For the negative pulse $Z\leq0$, both horizontal components are
strictly negative throughout this domain. Their vorticity equations
then have strictly positive temperature forcing while still retaining
their old vorticity data and their damping. In particular the small
negative transition value of $\Omega_2/\sigma_2$ is not a bound
to carry unchanged through steering. Its full coupled growth must be
estimated from \eqref{tr:closed} before an endpoint sign is claimed.

There is an exact invariant return submanifold for the holding
equations: prescribe $Z=0$ and set $\Omega_3=0$. The vorticity
row in \eqref{tr:physical_pair} then vanishes identically, while
\begin{equation}\label{tr:hold}
 \begin{gathered}
 \zeta_3=(0,\sqrt R),\qquad
 \zeta_2=\frac{(-b,g)}{\sqrt R},\qquad
 \zeta_1=\frac{(-y,x)}{\sqrt R},\\
 \dot\Omega_1=\frac{E_1y}{\sqrt R}-dK_1^2\Omega_1,\qquad
 \dot\Omega_2=\frac{E_2b}{\sqrt R}-dK_2^2D\Omega_2,\\
 \dot\Theta_3=-dK_3^2R\Theta_3,\qquad
 \dot\alpha=-\frac{\Omega_1y^2}{R}
                 -\frac{\Omega_2b^2}{DR}.
 \end{gathered}
\end{equation}
The $h,g,E_1,E_2$ rows remain those of \eqref{tr:closed}, and
the two numerator loss rows remain \eqref{tr:losses}. This
restriction of the full vector field proves invariance of the
return submanifold, and exhibits the continued motion of both
parents. It does not assert that a source trial has reached it.
The total local velocity curl on this submanifold is exactly
\begin{equation}\label{tr:origin_curl}
 2\dot\alpha+\Omega_1+\Omega_2
 =\Omega_1\left(1-\frac{2y^2}{R}\right)
  +\Omega_2\left(1-\frac{2b^2}{DR}\right).
\end{equation}
This follows by taking the curl of the rotational matrix and of
each rank-one shear; their respective curls are $2\dot\alpha$,
$\Omega_1$, and $\Omega_2$. No sign for the full third-return
curl follows without proving that endpoint's parent amplitudes.

Finally the exact temperature gain and diffusion exponents on any
constructed finite quotient interval $[0,\tau_f]$ are
\begin{equation}\label{tr:costs}
 \begin{gathered}
 \log\frac{P_3(\tau_f)}{P_{3,1}}
 =\int_0^{\tau_f}kv\,d\tau
          -\delta_3\int_0^{\tau_f}R\,d\tau,\\
 \delta_1\tau_f,\qquad
 \delta_2\int_0^{\tau_f}D\,d\tau,\qquad
 \delta_3\int_0^{\tau_f}R\,d\tau.
 \end{gathered}
\end{equation}
Each is exactly $dK_j^2\int|\zeta_j(t)|^2dt$ in original
physical time. The two explicit losses in \eqref{tr:losses}
also remain part of the older response; common damping of the
newest quotient has not removed them. Compact original-profile
PDE realization additionally uses the scalar and full vector
localization forces already calculated from the exact affine
background. Nothing in this steering coordinate map discards
those spatial forces.
